\documentclass[addpoints]{exam}
%\documentclass[9pt]{extarticle}
\pagestyle{headandfoot}

\usepackage[T1,T2A]{fontenc}
\usepackage[utf8]{inputenc}
%\usepackage[russian]{babel} % changes date and table caption

\usepackage{amsmath, amsfonts, amssymb}
\usepackage{verbatim}
\usepackage{graphicx}
\usepackage[usenames,dvipsnames]{xcolor}
\usepackage{parskip}
\usepackage{tikz} % drawing schemes
\usetikzlibrary{arrows.meta, positioning, calc, fit, backgrounds, shapes.geometric}

% fancy table
\usepackage{caption}
\usepackage{tabularray}
\usepackage{varwidth}
\UseTblrLibrary{booktabs}
\setlength{\heavyrulewidth}{1.5pt}
\usepackage{enumitem}
\usepackage{subcaption}
\usepackage{multirow}
\usepackage{microtype}
\usepackage{makecell}
%\usepackage{etoolbox} % For logic

% ---------- Python code listings ----------
\usepackage{listings}
\definecolor{codekeyword}{RGB}{0,0,180}
\definecolor{codestring}{RGB}{163,21,21}
\definecolor{codecomment}{RGB}{0,128,0}
\definecolor{codeident}{RGB}{38,127,153}
\definecolor{codenumber}{RGB}{125,125,125}
\definecolor{codeback}{RGB}{248,248,248}
\definecolor{codeframe}{RGB}{220,220,220}
\lstdefinestyle{python}{
  language=Python,
  basicstyle=\ttfamily\small,
  keywordstyle=\color{codekeyword}\bfseries,
  stringstyle=\color{codestring},
  commentstyle=\color{codecomment}\itshape,
  emph={self},emphstyle=\color{codeident}\itshape,
  showstringspaces=false,
  breaklines=true,
  frame=single,
  framesep=5pt,
  rulecolor=\color{codeframe},
  backgroundcolor=\color{codeback},
  numbers=left,
  numberstyle=\ttfamily\footnotesize\color{codenumber},
  numbersep=8pt,
  stepnumber=1,
  xleftmargin=22pt,
  xrightmargin=4pt,
  aboveskip=8pt,
  belowskip=4pt,
  columns=fullflexible,
  keepspaces=true,
}
\lstset{style=python}
% ------------------------------------------

% ---------- float placement ----------
% Keep every figure inside the section it belongs to, and give LaTeX
% enough slack on a page that it stops deferring floats to the very end.
\usepackage[section]{placeins}
\setcounter{topnumber}{3}
\setcounter{bottomnumber}{2}
\setcounter{totalnumber}{5}
\renewcommand{\topfraction}{0.9}
\renewcommand{\bottomfraction}{0.7}
\renewcommand{\textfraction}{0.08}
\renewcommand{\floatpagefraction}{0.75}
% -------------------------------------

\DeclareMathOperator*{\argmax}{arg\,max}
\DeclareMathOperator*{\argmin}{arg\,min}

\newcommand{\semester}{SS 2026}
\runningheader{\students}{}{\semester}
\runningfooter{}{\thepage}{}
\headrule
\footrule

% ---------- Modify team name, students (matriculation number), exercise number here ----------
\newcommand{\teamname}{adl2026-team1}
\newcommand{\students}{Velikanov Maksim}
\newcommand{\assignmentnumber}{1}
% ---------- End Modify ----------

\title{Leveraging world model internal representation for ego trajectory (?) prediction with an attentive diffusion probe}
\author{\students}
\date{\today}

\begin{document}
    \maketitle

    \begin{abstract}
    This report studies what the internal representation of the Orbis driving
    world model contains. We train a small attentive diffusion probe on the
    extracted image latents and use it to predict the future 2D trajectory of the
    vehicle. Most of the report is the background needed for that study. We
    describe the two-stage Orbis model, a hybrid VAE tokenizer followed by a
    flow-matching latent world model, and the probing setup built on top of it.
    We then review diffusion and flow matching, the diffusion-transformer
    architecture behind current generative models, world models for driving and
    navigation, and the trajectory-planning literature that produces the same
    kind of output as our probe. The data and metrics come from the nuPlan
    dataset and the NAVSIM pseudo-simulation benchmark with its PDM/EPDMS
    scoring. The last two sections report the experiments run so far: a
    controlled synthetic study that checks the flow-matching denoiser against a
    trajectory distribution known in closed form, and the probe as it is actually
    implemented on real NAVSIM driving logs, including the whole-episode padded
    training setup and the multi-token image aggregator through which the frozen
    latent conditions the denoiser.
    \end{abstract}

    \tableofcontents
    \clearpage

    \section{Introduction}
\label{sec:introduction}

A lot of recent work on autonomous driving is organized around the idea that an
agent should be able to imagine what will happen before it acts. A \emph{world
model} is the component that makes that possible: it learns the dynamics of the
environment from observation and can roll out plausible future states given the
past. When the observations are front-camera video, what comes out is a learned
simulator of the road ahead.

Two things make this attractive. Raw driving video exists in enormous quantities
and needs no manual annotation, so a world model can in principle pick up the
long tail of driving behaviour that would be expensive to label. And a model
that predicts future frames has to learn something about perception, about how
the scene moves, and about the ego vehicle's own motion, which suggests its
internal representation could be useful for decisions further down the stack.

This report takes one such model, \textbf{Orbis}~\cite{orbis}, and asks a narrow
question about it: does the latent it produces already contain what is needed to
predict where the car is going? We answer it by attaching a small
\emph{attentive diffusion probe} to the frozen image latent and training it to
predict the future 2D trajectory of the ego vehicle. The probe is kept small on
purpose. If it works, we want to be able to say the information was already in
the representation, rather than that the probe learned the task on its own. We
look at pointwise accuracy and also at the \emph{distribution} of the predicted
trajectories, since driving is uncertain and multi-modal, and a model that
collapses onto a single averaged path is not capturing that.

The rest of the report is background for this study. The introduction describes
Orbis and the probing setup. Section~\ref{sec:related} covers diffusion and flow
matching, the diffusion transformer, world models, and the planning literature.
Section~\ref{sec:datasets} describes nuPlan and the NAVSIM pseudo-simulation
benchmark, which is where trajectories get produced and scored.
Section~\ref{sec:synthetic} checks the generative core of the probe on a
synthetic distribution whose ground truth we know in closed form, and
Section~\ref{sec:navsim-probe} describes the probe as it is actually implemented
on NAVSIM logs: the data pipeline, the variable-length training setup, the image
conditioning, and the metrics we use to judge a distribution of futures.

\subsection{The Orbis world model}
\label{sec:orbis}

Orbis~\cite{orbis} predicts future video frames from front-camera observations.
Its main claim is that a continuous, flow-matching world model is cheaper and
holds up better over long rollouts than a discrete, token-based one. With $469$M
parameters and around $280$ hours of raw video it reaches state-of-the-art
video-prediction quality on the nuPlan, Waymo and nuPlan-turns benchmarks, and
stays stable on rollouts of up to twenty seconds. For our purposes the important
detail is that it is trained on raw video and nothing else. No depth maps, no
bounding boxes, no HD maps, no privileged supervision of any kind, so whatever
ends up in the latent got there through the predictive objective.

The architecture follows the two-stage recipe that is by now standard for latent
generative models, shown in Figure~\ref{fig:orbis}. \textbf{Stage~1} is an image
\emph{tokenizer}, an autoencoder in the spirit of a VAE, which compresses a
frame into a compact latent and decodes it back to pixels. \textbf{Stage~2} is
the latent \emph{world model}, which works entirely in that latent space and
autoregressively generates the next frame's latent from a window of context
frames. Splitting the model this way lets the expensive generative part run at
low spatial resolution while the tokenizer handles the high-frequency pixel
detail.

\paragraph{Stage 1: hybrid image tokenizer.}
The tokenizer maps an image $\mathbf{I}\in\mathbb{R}^{H\times W\times 3}$ to a
spatial latent compressed by a factor of $16$ in each dimension. What makes it
unusual is that it is \emph{hybrid} and \emph{factorized}. Two encoders produce
complementary representations: a \emph{semantic} encoder $\mathcal{E}_s$, kept
close to DINOv2 features by a cosine similarity loss, and a \emph{detail}
encoder $\mathcal{E}_d$ that picks up low-level appearance. Each stream is
quantized against its own codebook, giving semantic tokens $\mathbf{q}_s$ and
detail tokens $\mathbf{q}_d$, and the tokenizer is trained with the VQGAN
objective. During fine-tuning the vector-quantization bottleneck is bypassed
with probability $50\%$, so the same encoder ends up supporting a continuous
latent $\mathbf{x}=(\mathbf{x}_s;\mathbf{x}_d)$ as well as a discrete one
$\mathbf{q}=(\mathbf{q}_s;\mathbf{q}_d)$. That is what lets Orbis compare
continuous and discrete world models under identical conditions. The latent we
probe in this report is the continuous one, $\mathbf{x}$.

\paragraph{Stage 2: latent-space diffusion world model.}
The world model is a next-frame autoregressive predictor. It takes a window of
encoded context frames $\mathbf{x}_{0:N-1}$ together with a noised (or fully
masked) target frame, and predicts the clean target latent $\mathbf{x}_N$. In
the continuous case this is done with a flow-matching objective and a
spatio-temporal diffusion transformer (STDiT). Each block applies spatial
attention within a frame, temporal attention across frames, and a feed-forward
layer, and conditioning signals come in through adaptive layer normalization
(AdaLN). At inference the model draws a noise tensor for the target frame and
integrates the learned velocity field back to a clean latent. The generated
frame is then appended to the context and the oldest one is dropped, so the
window slides forward and the rollout continues. Ego-motion control (steering
angle and speed) is optional: a small MLP embeds it and adds it to the AdaLN
conditioning, which shows that the latent dynamics can be steered.

The flow-matching objective defines a forward path from data to noise by linear
interpolation,
\begin{equation}
\mathbf{x}^{\tau} = (1-\tau)\,\mathbf{x} + \tau\,\boldsymbol{\epsilon},
\qquad \tau\in[0,1],\quad \boldsymbol{\epsilon}\sim\mathcal{N}(0,I),
\label{eq:fm-forward-intro}
\end{equation}
and trains a network to predict the velocity that transports the corrupted
target frame $\mathbf{x}_N^{\tau}$ back towards the data manifold, conditioned
on the context:
\begin{equation}
\mathcal{L} = \mathbb{E}_{\tau,\boldsymbol{\epsilon}}
\big[\,\lVert \mathbf{v}(\mathbf{x}_N^{\tau};\mathbf{x}_{0:N-1})
- (\boldsymbol{\epsilon}-\mathbf{x}_N) \rVert^2\,\big].
\label{eq:fm-loss-intro}
\end{equation}
Section~\ref{sec:diffusion-theory} goes through the theory behind
Equations~\eqref{eq:fm-forward-intro} and~\eqref{eq:fm-loss-intro}.

\begin{figure}[htbp]
\centering
\includegraphics[width=\textwidth,height=0.42\textheight,keepaspectratio]{../diagrams/diagram-orbis.pdf}
\caption{The two-stage Orbis architecture. \textbf{Stage~1} is a hybrid
tokenizer with a DINOv2-regularized semantic encoder and a detail encoder,
producing a continuous latent $\mathbf{x}$, or equivalently a set of discrete
tokens. \textbf{Stage~2} is a spatio-temporal diffusion transformer trained with
flow matching, which predicts the next-frame latent $\hat{\mathbf{x}}_N$ from a
window of context latents. The prediction is appended and the window slides
forward, which is how long rollouts are produced. The frozen latent
$\mathbf{x}$ is the representation we probe.}
\label{fig:orbis}
\end{figure}

\subsection{Probing the latent for future trajectories}
\label{sec:probe}

The usual way to find out what a learned representation contains is to freeze it
and train a small read-out head on some target task. If a simple probe succeeds,
the information is accessible in the representation without much further
processing~\cite{dosovitskiy_vit}. Our target task is future trajectory
prediction: given the latent of the current and recent frames, predict the
sequence of future 2D positions $(x,y)$ that the ego vehicle will follow in its
own coordinate frame. This is the same output space that end-to-end planners use,
and the same one the trajectory-based evaluation in Orbis uses, so the
predictions can be compared against driving benchmarks directly.

Two choices define the probe. The first is to use an \emph{attention} probe
instead of a plain linear layer. The image latent is a spatial grid of tokens,
and which regions matter for the future trajectory depends on the scene: lane
markings, the vanishing point, a vehicle in front. A small set of query tokens
attends over the frozen grid and pools it into a fixed-size context vector, so
the probe can concentrate on the informative regions without touching the
backbone. The second is to produce the trajectory with a \emph{diffusion} model
operating directly in 2D trajectory space. Driving is multi-modal. At an
intersection, turning left and turning right are both correct, and a regression
head trained with an $L_2$ loss would average them into a straight-ahead path
that no driver would take. A diffusion head learns the conditional distribution
instead, and can sample several distinct plausible futures, which is what we
need if we want to say anything about the \emph{distribution} of behaviours
stored in the latent. Figure~\ref{fig:probe} sketches the setup. The version in
Section~\ref{sec:navsim-probe} departs from the sketch in one place: the query
tokens are fixed sinusoidal slot embeddings, and the pooled tokens are not
collapsed into a single conditioning vector but enter the denoiser as keys and
values of its attention stack.

\begin{figure}[htbp]
\centering
\includegraphics[width=\textwidth,height=0.40\textheight,keepaspectratio]{../diagrams/diagram-probe.pdf}
\caption{The attentive diffusion probe. The Orbis encoder is frozen and produces
a spatial latent grid for the recent frames. A small attention probe with query
tokens pools that grid into a context vector $\mathbf{c}$, which conditions a
diffusion model that generates future 2D ego trajectories $(x,y)$ directly in
trajectory space. Because the head is generative it can sample several distinct
futures, for instance going straight and turning, so we can evaluate the
distribution of trajectories rather than a single averaged path.}
\label{fig:probe}
\end{figure}

This setup touches most of the report. The generative machinery in both Orbis
and the probe is diffusion and flow matching (Section~\ref{sec:diffusion-theory}),
carried by the diffusion transformer of Section~\ref{sec:dit}. Reusing a
predictive video model as a substrate for planning is what places the work in
the world-model literature of Section~\ref{sec:world-models}, next to navigation
world models that already use rollouts to make decisions, and in the planning
literature of Section~\ref{sec:planners}. Evaluating the trajectories that come
out relies on the datasets and metrics of Section~\ref{sec:datasets}.

    \section{Background and Related Work}
\label{sec:related}

\subsection{The Orbis world model}
\label{sec:orbis}

Orbis~\cite{orbis} predicts future front-camera frames from past ones. It is
trained on about 280 hours of raw driving video, with no labels, maps or depth.
Whatever its features know about driving was learned from the prediction
objective alone. The model has two stages (Figure~\ref{fig:orbis}).

The first stage is an image tokenizer. It compresses a frame by a factor of 16 in
each dimension into two latent maps: a semantic latent $q_{\text{sem}}$, trained
to stay close to DINOv2 features, and a detail latent $q_{\text{rec}}$ for the
appearance that the decoder needs. Both have 16 channels.

The second stage is the world model itself, a spatio-temporal diffusion
transformer with 24 blocks and about 470M parameters. It takes the latents of
the context frames plus a noisy latent for the next frame, and predicts the
flow-matching velocity that denoises the next frame. Each block mixes tokens
within a frame and across frames. We probe both stages: the tokenizer latents and
the tokens inside the world model after a given block.

\begin{figure}[htbp]
\centering
\includegraphics[width=\textwidth,height=0.36\textheight,keepaspectratio]{../diagrams/diagram-orbis.pdf}
\caption{The two-stage Orbis model. Stage 1 is the tokenizer with a
DINOv2-regularized semantic branch and a detail branch. Stage 2 is a
spatio-temporal diffusion transformer that predicts the next-frame latent
$\hat{\mathbf{x}}_N$ from the context latents. We read features from both
stages.}
\label{fig:orbis}
\end{figure}

\subsection{Flow matching}
\label{sec:flow-matching}

Both Orbis and our probe are flow-matching models~\cite{lipman_fm}, a close
relative of diffusion models~\cite{ho_ddpm}. A sample $\mathbf{x}$ from the data
is mixed with Gaussian noise $\boldsymbol{\epsilon}$ along a straight line,
\begin{equation}
\mathbf{x}^{\tau} = (1-\tau)\,\mathbf{x} + \tau\,\boldsymbol{\epsilon},
\qquad \tau\in[0,1],
\label{eq:fm-path}
\end{equation}
and a network $\mathbf{v}_\theta$ learns the velocity of this path given a
condition $\mathbf{c}$:
\begin{equation}
\mathcal{L} = \mathbb{E}_{\mathbf{x},\mathbf{c},\tau,\boldsymbol{\epsilon}}
\big\lVert \mathbf{v}_\theta(\mathbf{x}^{\tau},\tau,\mathbf{c})
- (\boldsymbol{\epsilon}-\mathbf{x}) \big\rVert^2 .
\label{eq:fm-loss}
\end{equation}
To sample, we draw $\boldsymbol{\epsilon}$ and integrate
$d\mathbf{x}^\tau/d\tau=\mathbf{v}_\theta$ from $\tau=1$ to $\tau=0$ with a few
Euler steps.

One property matters later. At $\tau=1$ the input is pure noise, so the best
the network can predict is
$\mathbf{v}^\ast(\boldsymbol{\epsilon},1,\mathbf{c}) =
\boldsymbol{\epsilon}-\mathbb{E}[\mathbf{x}\mid\mathbf{c}]$. A single Euler step
from $\tau=1$ to $0$ therefore lands every sample on the conditional mean
$\mathbb{E}[\mathbf{x}\mid\mathbf{c}]$, whatever the noise. With few steps a
flow model looks unimodal even if it has learned a multimodal distribution.

\subsection{Memorization in diffusion models}
\label{sec:memorization}

A diffusion model trained to its global optimum on a finite training set
reproduces the training samples. Real models usually generalize anyway. Bonnaire
et al.~\cite{bonnaire} explain this with two time scales: the model reaches good
samples after a training time $\tau_{\text{gen}}$ that does not depend on the
dataset size $n$, and starts to memorize after $\tau_{\text{mem}}\propto n$.
Between the two there is a window of generalization that grows with the data.
Favero et al.~\cite{favero} observe the same linear law across image and
language diffusion models and propose early stopping that scales with $n$. Our
probe is a conditional diffusion model on a small target, so we can test this law
directly on driving data (Section~\ref{sec:scaling}).

\subsection{Probing frozen representations}
\label{sec:probing}

A probe is a small model trained on frozen features to read out a target
quantity~\cite{alain_probe}. A linear probe tests what is linearly available. A
small nonlinear probe tests what is available with little extra computation. We
use both. We compare Orbis with two strong self-supervised encoders.
DINOv2~\cite{dinov2} is an image encoder trained by self-distillation, and the
Orbis semantic latent is trained to imitate it. V-JEPA~2~\cite{vjepa2} is a video
encoder trained to predict masked parts of a clip in feature space, and was also
used for planning in robotics. Both are trained on general data, Orbis on driving
video only.

\subsection{Planning on NAVSIM}
\label{sec:planners}

NAVSIM~\cite{navsim1,navsim} scores a planner's 4-second trajectory by
simulating it in the logged scene (Section~\ref{sec:epdms}). Most strong
planners on it are trained by imitation of the logged driver.
TransFuser~\cite{transfuser} regresses one trajectory from camera and LiDAR.
VADv2~\cite{vadv2} and Hydra-MDP~\cite{hydra_mdp,hydra} score a fixed vocabulary
of candidate trajectories, and Hydra-MDP adds heads that learn the simulator's
sub-scores for every candidate. DiffusionDrive~\cite{diffusiondrive} generates
trajectories with a truncated diffusion process around anchor trajectories.
DriveSuprim~\cite{drivesuprim} and similar methods focus on selecting the best
of many candidates. PDM-Closed~\cite{pdm} is a rule-based planner with access to
the map and the other agents, and is the reference for the progress score.

Two warnings from this literature shape our setup. First, the ego status
(velocity and acceleration) alone already predicts the future trajectory well in
open-loop benchmarks~\cite{admlp,egostatus}. A probe that gets the ego status
measures something different from a probe that sees only images, so we report
the two separately. Second, average displacement error is a poor guide to
driving quality, which is why we also report EPDMS.

World models have been used for planning before. Drive-WM~\cite{drivewm} rolls
out a world model to compare candidate plans, LAW~\cite{law} and
World4Drive~\cite{world4drive} train a latent world model as an auxiliary task of
an end-to-end planner, and navigation world models~\cite{nwm} plan by search
over rollouts. These methods train the world model together with the planner or
use it as a simulator. We instead ask what a world model trained without any
planning objective already knows, by keeping it frozen.

    \section{Data and Evaluation Protocol}
\label{sec:data}

\subsection{NAVSIM scenes}
\label{sec:navsim}

NAVSIM~\cite{navsim1} is built on the OpenScene version of nuPlan~\cite{nuplan},
real driving logs from four cities: Las Vegas, Boston, Pittsburgh and Singapore.
Singapore drives on the left. A NAVSIM scene is a window of a log sampled at
2\,Hz. The input is four history frames covering 1.5\,s, together with the ego
velocity, acceleration and a driving command (left, straight or right at the next
junction, taken from the route). The target is the ego trajectory over the next
4\,s, eight poses $(x,y,\psi)$ in the ego frame of the current scene. We use only
the front camera, since that is what Orbis was trained on.

\paragraph{Why we changed the protocol.}
Our earlier probe used its own data pipeline. It cut each log into the longest
run of frames with images on disk, which gave 9,577 episodes, and predicted the
trajectory for up to 20\,s. It also fed the image at $t=0.5$\,s as an input while
the motion up to $0.5$\,s was part of the target, which leaks the answer.
Following the official protocol fixes these problems, gives 10.7 times more
training scenes (the scenes overlap, but each has its own target), and makes our
EPDMS comparable with published planners.

\paragraph{Splits.}
Table~\ref{tab:splits} lists the splits. Consecutive navtrain scenes share 13 of
their 14 frames, so a random split by scene would put near-copies of the
validation scenes into training. We hold out whole logs instead, 10\% of them,
chosen by a hash of the log name. All model selection uses these held-out logs.
navtest and navhard come from separate test logs and are used only for the final
numbers. In navhard the second stage starts from synthetic scenes, rendered with
Gaussian splatting from an ego pose displaced from the logged one, so it tests
whether a planner recovers from a state the driver never was in.

\begin{table}[htbp]
\centering
\caption{Data splits. The validation split is 98 navtrain logs held out from
training. navtest and navhard come from the OpenScene test logs.}
\label{tab:splits}
\begin{tblr}{colspec={lllr}, row{1}={font=\bfseries}}
\toprule
Split & Role & Source & Scenes \\
\midrule
train      & training                    & navtrain, 1,094 logs & 92,381 \\
val        & model selection             & navtrain, 98 held-out logs & 9,675 \\
navtest    & test, one stage             & test logs & 12,146 \\
navhard    & hard test, two stages       & test logs, partly synthetic & 5,915 \\
\bottomrule
\end{tblr}
\end{table}

\subsection{EPDMS}
\label{sec:epdms}

NAVSIM scores a trajectory by simulating it for 4\,s in the logged scene, with
the other agents replaying their logged motion. The Extended Predictive Driver
Model Score~\cite{navsim,hydra} multiplies hard penalties with a weighted average
of soft terms. In the devkit version we use,
\begin{equation}
\text{EPDMS} =
\text{NC}\cdot\text{DAC}\cdot\text{DDC}\cdot\text{TLC}\cdot
\frac{5\,\text{EP} + 5\,\text{TTC} + 2\,\text{LK} + 2\,\text{HC} + 2\,\text{EC}}{16}.
\label{eq:epdms}
\end{equation}
The penalties are no at-fault collision (NC), drivable-area compliance (DAC),
driving-direction compliance (DDC) and traffic-light compliance (TLC). The soft
terms are ego progress relative to PDM-Closed (EP), time to collision (TTC), lane
keeping (LK), history comfort (HC) and two-frame extended comfort (EC). EC
compares the plan with the plan made 0.5\,s earlier, so a sampler that jumps
between modes from frame to frame loses it. A penalty is waived if the human
driver made the same violation. On navhard the scores of the two stages are
combined, with stage-two scenes weighted by how close their start is to the
endpoint of the stage-one plan.

A full navtest run takes about 20 minutes. All EPDMS numbers in this report come
from the official devkit. Our devkit gives the logged human trajectory 0.949 and
PDM-Closed 0.871 on a 500-scene navtest subset, a few points above the published
0.903 and 0.828. Comparisons with published navtest numbers are therefore
approximate. Our navhard numbers are not comparable with the public leaderboard
at all, which runs on a private copy of the split: constant velocity scores 0.326
in our run and 0.114 on the leaderboard. On navhard we only compare our own runs.

    \section{A Controlled Synthetic Study}
\label{sec:synthetic}

Before putting the diffusion probe on the high-dimensional Orbis latent and on
noisy real driving logs, we test its generative core, meaning the flow-matching
objective together with the temporal-transformer denoiser, on a controlled
synthetic distribution of 2D trajectories whose shape we know in closed form.
The reason is fault isolation. If the model fails to reproduce the trajectory
distribution on real nuPlan data, the cause could be the data pipeline, the
tokenizer, the normalization, the conditioning or the diffusion model itself,
and telling those apart after the fact is hard. Replacing the conditioning
latent with a synthetic signal and the target with an analytically generated
curve strips the system down to its generative part and leaves one question:
\emph{can a flow-matching temporal transformer denoise a whole trajectory at
once and reproduce the target distribution as a set?} The module tested here
(\texttt{DiffusionModel}) is the same one we later use on real data, apart from
the added image conditioning, so whatever we conclude here carries over.

The task is kept as small as possible. Each trajectory is a smooth cubic
B\'ezier curve starting at the origin, leaving along $+x$, and fanning out left
or right depending on one scalar \emph{turn} parameter. The model has to denoise
the whole trajectory jointly rather than autoregressively, under a flow-matching
objective, and the samples it produces should reproduce the fan-shaped
ground-truth distribution \emph{as a set}. Section~\ref{sec:syn-conditioning}
explains that the turn conditioning is switched off on purpose in this
experiment, so what we measure is \emph{distribution coverage} and not
per-sample turn accuracy.

\subsection{Synthetic data generation}
\label{sec:syn-data}

The dataset is procedural. Nothing is stored on disk; every sample is computed
on the fly from a turn parameter drawn uniformly,
$\text{turn}\sim\mathcal{U}[0,1]$. The turn maps to a heading angle and to an
endpoint on a circle of radius $r$,
\begin{equation}
\alpha(\text{turn}) = 30^\circ - 60^\circ\cdot\text{turn},
\qquad
(X, Y) = \big(r\cos\alpha,\; r\sin\alpha\big),
\label{eq:syn-endpoint}
\end{equation}
so $\text{turn}=0$ gives a left turn at $+30^\circ$, $\text{turn}=0.5$ goes
straight ahead, and $\text{turn}=1$ gives a right turn at $-30^\circ$. The
configuration uses $r=1$, so all endpoints lie on the unit circle inside a
$60^\circ$ wedge centred on the $+x$ axis. That wedge is the fan visible in the
``ground truth'' panels of every result figure in
Section~\ref{sec:syn-ablations}.

Given the start $p_0=(0,0)$ and the endpoint $p_3=(X,Y)$, the curve is a cubic
B\'ezier
\begin{equation}
B(s) = (1-s)^3\,p_0 + 3(1-s)^2 s\,p_1 + 3(1-s)\,s^2\,p_2 + s^3\,p_3,
\qquad s\in[0,1],
\label{eq:bezier}
\end{equation}
with both interior control points placed so that the curve starts out straight,
$p_1 = p_2 = (0.5\,r,\,0)$. Since they both sit on the $+x$ axis, the tangent at
the start points along $+x$, and the curve only bends towards its endpoint in its
second half. That is why every ground-truth trajectory leaves the origin along
the same straight stem before fanning out, the shared horizontal segment around
$x\in[0,0.2]$ that is easy to see in the plots. The curve is sampled at $T$
equally spaced parameter values $s\in\{0,\tfrac{1}{T-1},\dots,1\}$, which gives a
$(T,2)$ array of positions. The configuration uses $T=20$.

\paragraph{Velocity representation.}
The model does not operate on absolute positions but on \emph{velocities}, i.e.\
first-order position deltas,
\begin{equation}
v_t = p_{t+1} - p_t,\qquad t=0,\dots,T-2,
\label{eq:velocity}
\end{equation}
with the last entry zero-padded since there is no frame after the last:
\begin{lstlisting}
velocity = np.zeros_like(trajectory)
velocity[:-1] = trajectory[1:] - trajectory[:-1]   # (T, 2)
\end{lstlisting}
A sample is therefore a triple $(\text{turn},\,\text{position}\in\mathbb{R}^{T\times2},\,
\text{velocity}\in\mathbb{R}^{T\times2})$. Velocity space is easier to work in
for two reasons. The signal is roughly zero-mean and stationary along the
trajectory, which suits a diffusion model better than absolute coordinates that
grow monotonically in $x$. And positions come back with a cumulative sum at
sampling time (Section~\ref{sec:syn-flow}). The price is that per-step velocity
errors accumulate in that sum, so a small systematic bias in the predicted deltas
shows up as endpoint drift rather than as a wobbly path. This turns out to
matter: endpoint scatter ends up being the most sensitive quality indicator we
have. Figure~\ref{fig:syn-pipeline} summarizes the generation and representation
pipeline, and Table~\ref{tab:syn-config} lists the configuration.

\begin{figure}[htbp]
\centering
\includegraphics[width=\textwidth,height=0.34\textheight,keepaspectratio]{../diagrams/diagram-syn-pipeline.pdf}
\caption{Synthetic data generation and representation. A scalar turn parameter
fixes an endpoint on a unit arc, a cubic B\'ezier with straight-stem control
points interpolates from the origin to that endpoint, and the sampled positions
are differenced into velocities. Those are standardized per channel
(Section~\ref{sec:syn-norm}) before they reach the diffusion model. At sampling
time positions are reconstructed from velocities by cumulative summation.}
\label{fig:syn-pipeline}
\end{figure}

\begin{table}[htbp]
\centering
\caption{Synthetic dataset and training configuration. The image-related fields
inherited from the nuPlan configuration (HDF5 paths, augmentation, frame rates)
are ignored by the synthetic loader, so they are left out here.}
\label{tab:syn-config}
\begin{tblr}{
  colspec={llX[l]},
  row{1}={font=\bfseries},
}
\toprule
Parameter & Value & Role \\
\midrule
$T$            & $20$    & timesteps per trajectory \\
$r$            & $1$     & endpoint radius (unit arc) \\
turn range     & $[0,1]$ & full left-to-right fan ($+30^\circ$ to $-30^\circ$) \\
dataset size   & $10000$ & trajectories per epoch \\
batch size     & $30$    & per-GPU batch \\
base LR        & $10^{-3}$ & AdamW, weight decay $0.01$, cosine decay to $0.1\times$ \\
precision      & 16-mixed & with gradient clipping (norm $1.0$) \\
training steps & $5000$  & for every ablation run \\
diffusion steps & $20$   & Euler steps at sampling time \\
$\sigma_{\min}$ & $10^{-5}$ & flow-matching schedule floor \\
\bottomrule
\end{tblr}
\end{table}

\subsection{Velocity normalization}
\label{sec:syn-norm}

The two velocity channels have very different statistics. The $x$-deltas are
always positive and fairly large, because the trajectory marches steadily in
$+x$. The $y$-deltas average out near zero and are symmetric, because left turns
cancel right turns across the dataset. A diffusion model injects
\emph{isotropic} Gaussian noise, so unequal channel scales make the effective
noise level depend on the channel: one shared noise magnitude is too large for
the small-scale channel and too small for the other one at the same time.
Standardizing each channel to zero mean and unit variance gets rid of the
mismatch.

The statistics are computed once over the whole dataset, per channel, and cached
to disk so that repeated runs with the same $(\text{dataset size},T,r)$ skip the
recomputation:
\begin{equation}
\mu = \frac{1}{NT}\sum_{n,t} v^{(n)}_t \in \mathbb{R}^2,
\qquad
\sigma = \Big[\tfrac{1}{NT}\sum_{n,t}\big(v^{(n)}_t-\mu\big)^2\Big]^{1/2}
\in \mathbb{R}^2.
\label{eq:syn-stats}
\end{equation}
With a small constant $\varepsilon=10^{-8}$ for numerical safety, the forward
transform applied per channel in the data loader and its exact inverse applied to
the model's sampled velocities are
\begin{equation}
\hat v = \frac{v-\mu}{\sigma+\varepsilon},
\qquad\qquad
v = \hat v\,(\sigma+\varepsilon) + \mu.
\label{eq:syn-norm}
\end{equation}
The same $(\sigma+\varepsilon)$ factor appears in both directions, so the round
trip is exact up to floating point. Since the model only ever sees standardized
velocities during training, the diffusion process runs in a space where both
channels are approximately $\mathcal{N}(0,1)$, which is the prior the sampler
starts from. The padded last velocity row, a zero, is included in the
statistics. With $T=20$ that is a $1/20$ contribution which shrinks $\mu$ and
$\sigma$ slightly. It does no harm, and more to the point it is identical on the
normalize and the denormalize side.

A boolean flag (\texttt{output\_normalization}) reduces both transforms in
Equation~\eqref{eq:syn-norm} to the identity. That toggle is what the
normalization ablation in Section~\ref{sec:syn-ablation-norm} switches. In the
production code path the cached statistics are replaced by an
exponential-moving-average module, which keeps the same per-channel
$(\mu,\sigma)$ and uses the same formulas but updates them online during training
instead of in one pass beforehand. The synthetic study uses the cached
dataset-side version, and the two are mathematically the same thing.

\subsection{Denoiser architecture}
\label{sec:syn-model}

The denoiser takes the noised velocity trajectory at diffusion time $t$ and
predicts a per-step velocity field. It has three inputs: the noisy deltas
$\hat v_t\in\mathbb{R}^{B\times T\times 2}$, a diffusion timestep
$t\in\mathbb{R}^B$ (one per trajectory, shared across the $T$ points), and a turn
scalar, which is disabled (Section~\ref{sec:syn-conditioning}). Two scalars are
lifted into embeddings of width $D$, the model's hidden dimension, by a
sinusoidal positional encoding followed by a two-layer MLP. One is a
\emph{diffusion-time} embedding $e_{\text{time}}\in\mathbb{R}^{B\times D}$
broadcast to all timesteps; the other is a \emph{frame-index} embedding
$e_{\text{frame}}\in\mathbb{R}^{T\times D}$ computed from the normalized index
$\text{arange}(T)/T$ and broadcast across the batch. The frame embedding is what
tells the model where along the trajectory each point sits. Attention is
permutation-equivariant, so without it the model would have no way to order the
points.

The per-point input is the concatenation of the noisy velocity, the two
embeddings, and the (zeroed) turn scalar,
\begin{equation}
x_i = \big[\,\underbrace{\hat v_i}_{2},\;
\underbrace{e_{\text{time}}}_{D},\;
\underbrace{e_{\text{frame},i}}_{D},\;
\underbrace{\text{turn}}_{1}\,\big]\in\mathbb{R}^{2+2D+1},
\label{eq:syn-input}
\end{equation}
reshaped to $(B\cdot T,\,2+2D+1)$ and pushed through a pointwise MLP of four
$\text{Linear}\to\text{Tanh}$ layers, which maps each point independently to a
$D$-dimensional feature. Nothing interacts across timesteps at this stage; it is
just a shared per-frame encoder. The temporal mixing happens in a stack of four
\texttt{TemporalBlock}s used in \emph{self-attention} mode, where the per-point
features act as query, key and value at once:
\begin{lstlisting}
pred_flow = features                       # (B, T, D)
for layer in self.cross_denoiser:
    pred_flow = layer(pred_flow, pred_flow, fm_emb)
\end{lstlisting}
Each block is a DiT-style transformer block~\cite{dit} with \emph{AdaLN-Zero}
conditioning. A LayerNorm without affine parameters is applied to the stream, and
the conditioning vector, which here is the diffusion-time embedding
$e_{\text{time}}$, goes through $\text{SiLU}\to\text{Linear}(D,8D)$ to give eight
$(B,D)$ modulation vectors: shift, scale and gate for the query, context and
output branches. The last linear is zero-initialized, so every block starts out
as an identity map and has to learn its way away from it. This is the AdaLN-Zero
initialization, and it is what makes deep diffusion transformers train stably.
An 8-head attention then lets each timestep gather information from the whole
trajectory, and the result is gated and added back residually. A feed-forward
sub-block ($\text{Linear}(D,2D)\to\text{GELU}\to\text{Linear}(2D,D)$) is
modulated, gated and added residually in the same way. The modulation helper
applies $x\mapsto x\,(1+\text{scale})+\text{shift}$. After the four blocks a
linear head $\text{Linear}(D,2)$ projects each point back to a 2D velocity, which
gives the predicted flow field $\hat u\in\mathbb{R}^{B\times T\times 2}$.
Figure~\ref{fig:syn-arch} shows the whole denoiser.

$D$ sets the width of the time and frame embeddings, the per-point MLP, the
attention embedding and the feed-forward dimension all at once, so it is the one
knob that controls capacity. That is what the first ablation in
Section~\ref{sec:syn-ablation-width} varies.

\begin{figure}[htbp]
\centering
\includegraphics[width=\textwidth,height=0.34\textheight,keepaspectratio]{../diagrams/diagram-syn-arch.pdf}
\caption{The flow-matching denoiser. A pointwise MLP encodes each timestep on its
own from the noisy velocity, the diffusion-time embedding and the frame-index
embedding. Four DiT-style temporal self-attention blocks, conditioned on the
diffusion time through AdaLN-Zero, then mix information across the $T$ timesteps,
and a linear head outputs the per-step flow-matching velocity field. The turn
conditioning is wired through but zeroed
(Section~\ref{sec:syn-conditioning}).}
\label{fig:syn-arch}
\end{figure}

\subsection{Flow-matching training and sampling}
\label{sec:syn-flow}

Training uses conditional flow matching~\cite{lipman_fm} with a linear,
rectified-flow-style interpolation schedule whose coefficients are
\begin{equation}
\alpha(t) = 1-t,
\qquad
\sigma(t) = \sigma_{\min} + t\,(1-\sigma_{\min}),
\qquad
\sigma_{\min} = 10^{-5}.
\label{eq:syn-schedule}
\end{equation}
Given clean normalized velocities $x=\hat v$ and noise
$\varepsilon\sim\mathcal{N}(0,I)$, the noised sample at diffusion time
$t\in[0,1)$ is the linear interpolant
\begin{equation}
x_t = \alpha(t)\,x + \sigma(t)\,\varepsilon
    = (1-t)\,x + \big(\sigma_{\min}+t(1-\sigma_{\min})\big)\,\varepsilon,
\label{eq:syn-forward}
\end{equation}
so $x_0 = x$ is clean data and $x_1\approx\varepsilon$ is essentially pure noise.
The diffusion time is drawn uniformly per trajectory, $t\sim\mathcal{U}[0,1)$.
The regression target is the time-derivative of the interpolation path, and for
the schedule in Equation~\eqref{eq:syn-schedule} it collapses to a
constant-coefficient combination of the data and the noise,
\begin{equation}
u_t = \dot\alpha(t)\,x + \dot\sigma(t)\,\varepsilon
    = x - (1-\sigma_{\min})\,\varepsilon,
\label{eq:syn-target}
\end{equation}
and the model is trained to regress it with a plain mean-squared error averaged
over all batch elements, timesteps, and channels,
\begin{equation}
\mathcal{L} = \big\lVert\, f_\theta(x_t,t) - u_t \,\big\rVert_2^2 .
\label{eq:syn-loss}
\end{equation}
The corresponding inner loop is compact:
\begin{lstlisting}
t = torch.rand(B, device=deltas.device)
noisy_deltas, noise = self.add_noise(deltas, t)          # noising x_t
pred_v = self.forward(turn, noisy_deltas, t)
target = self.A(t) * deltas + self.B(t) * noise          # target u_t
loss = ((pred_v - target) ** 2).mean()                   # MSE loss
\end{lstlisting}
The later image-conditioned model reserves a context prefix. Here the loss is
taken over all $T$ points.

\paragraph{Sampling.}
Generation integrates the learned ordinary differential equation
$\mathrm{d}x/\mathrm{d}t = f_\theta(x_t,t)$ backward from $t=1$ (noise) to $t=0$
(data) with a first-order Euler scheme over $20$ uniform steps:
\begin{lstlisting}
sampled = torch.randn(B, T, 2)                # x at t = 1
t_steps = torch.linspace(1, 0, num_steps + 1)
for i in range(num_steps):
    pred_v = self.forward(turn, sampled, t_steps[i].repeat(B))
    dt = t_steps[i] - t_steps[i + 1]          # > 0
    sampled = sampled + pred_v * dt
\end{lstlisting}
The integrated velocities are then denormalized by the inverse of
Equation~\eqref{eq:syn-norm} and integrated into positions by a cumulative sum
seeded at the origin (the synthetic experiment uses no ground-truth context
prefix),
\begin{equation}
\hat p_0 = 0,
\qquad
\hat p_{t} = \hat p_{t-1} + v_{t-1},
\qquad t=1,\dots,T-1.
\label{eq:syn-cumsum}
\end{equation}
The cumulative sum is where per-velocity errors compound. A small bias in
$\mathrm{d}x$ or $\mathrm{d}y$ builds up over the $20$-step integration, which is
why endpoint scatter rather than local path roughness is the most sensitive
indicator of sample quality in the figures below.

\subsection{Conditioning is disabled by design}
\label{sec:syn-conditioning}

One detail matters a lot for reading the results: the turn parameter is zeroed
inside the forward pass before it can reach the network. \texttt{turn} is wired
through the input concatenation of Equation~\eqref{eq:syn-input}, but it carries
no information, so the model is really an \emph{unconditional} generator of
trajectories. That has a few consequences. The model cannot map a particular
turn value to a particular endpoint, and it is not supposed to. The question to
ask when evaluating is whether the \emph{set} of generated trajectories covers
the same fan-shaped distribution as the ground truth. And in the figures the
colour of a generated curve is assigned by sample index, so it will not line up
with the same-coloured ground-truth curve; only the overall envelope should
match. Keeping the study unconditional focuses it on the property we need to
trust before scaling up, which is that the flow-matching transformer reproduces a
target trajectory \emph{distribution}.

\subsection{Ablations}
\label{sec:syn-ablations}

Everything below comes out of the same evaluation harness. It loads a trained
checkpoint, samples a batch of trajectories at turn values swept uniformly with
$\text{linspace}(0,1,N)$, denormalizes and integrates them, and plots the ground
truth next to the generated samples. Start points are circles, endpoints are
squares, the origin is a star, and colour encodes the swept turn index from blue
to red. As noted in Section~\ref{sec:syn-conditioning} the model is
unconditional, so the colours will not correspond between the two panels. What
matters is coverage of the fan.

\subsubsection{Model capacity}
\label{sec:syn-ablation-width}

The first ablation varies the capacity knob $D$, comparing $D=8$ with $D=16$ and
leaving everything else fixed (four temporal blocks, $5000$ training steps).
Figure~\ref{fig:syn-width} shows the result. At $D=8$ the model already recovers
the structure of the target: a shared straight stem out of the origin, a
left/right fan, and endpoints clustered near the unit arc around $x\approx0.9$ to
$1.0$. The fan is ragged. A few endpoints fall short of the arc and the angular
spacing of the generated curves is uneven, but the envelope is right. Doubling
the width to $D=16$ gives a cleaner and more evenly spread fan: the straight
initial stem is reproduced more faithfully, the individual curves are smoother,
and the endpoint cloud sits closer to the arc with fewer stragglers. What the
extra capacity mostly buys is smoother per-step velocities, and so less
cumulative-sum drift, plus better coverage of the extreme turns.

\begin{figure}[htbp]
\centering
\begin{subfigure}{\textwidth}
  \centering
  \includegraphics[width=0.95\linewidth]{../pictures/trajectory_comparison_8h_4layers_5000.png}
  \caption{$D=8$}
  \label{fig:syn-width-8}
\end{subfigure}

\vspace{0.8em}

\begin{subfigure}{\textwidth}
  \centering
  \includegraphics[width=0.95\linewidth]{../pictures/trajectory_comparison_16h_4layers_5000_steps.png}
  \caption{$D=16$}
  \label{fig:syn-width-16}
\end{subfigure}
\caption{Effect of the hidden width $D$, with four temporal blocks and $5000$
training steps. Ground truth is the left panel of each pair and the generated
samples the right one. The narrower $D=8$ model captures the fan but leaves a
ragged, unevenly spaced envelope with some short endpoints. The wider $D=16$
model gives smoother curves and a tighter endpoint cloud on the unit arc. Both
recover the correct distribution, so the task is learnable even at small
capacity, and $D=16$ is a reasonable default.}
\label{fig:syn-width}
\end{figure}

Capacity helps, but with diminishing returns. Even $D=8$ learns the task, and
$D=16$ is a good default for this problem. Table~\ref{tab:syn-ablation}
summarizes the comparison.

\subsubsection{Velocity normalization}
\label{sec:syn-ablation-norm}

The second ablation toggles \texttt{output\_normalization}, that is, whether the
transforms of Equation~\eqref{eq:syn-norm} are applied or replaced by the
identity. With normalization on the model trains and samples in standardized
velocity space; with it off the model sees raw velocities and the denormalization
step does nothing. Figure~\ref{fig:syn-norm} compares the two. With normalization
the generated trajectories are smooth, stay inside a sensible bounding box
($x\lesssim1.0$), and the endpoint cloud forms a recognizable arc covering the
full left-to-right fan. There is scatter, which is what we expect from an
unconditional model, but no gross scale errors and no runaway paths. Without
normalization the samples are much wavier. Several trajectories overshoot to
$x\approx1.2$, past the unit-radius ground truth, individual paths have kinks in
them instead of smooth arcs, and the endpoint cloud spreads out.

\begin{figure}[htbp]
\centering
\begin{subfigure}{\textwidth}
  \centering
  \includegraphics[width=0.95\linewidth]{../pictures/trajectory_comparison_with_norm.png}
  \caption{with normalization}
  \label{fig:syn-norm-on}
\end{subfigure}

\vspace{0.8em}

\begin{subfigure}{\textwidth}
  \centering
  \includegraphics[width=0.95\linewidth]{../pictures/trajectory_comparison_without_norm.png}
  \caption{without normalization}
  \label{fig:syn-norm-off}
\end{subfigure}
\caption{Effect of per-channel velocity normalization. With normalization (top)
the generated trajectories are smooth and bounded and the endpoint cloud is
arc-like. Without it (bottom) the samples are wavy and kinked, they overshoot to
$x\approx1.2$, and the endpoints scatter. The isotropic diffusion prior does not
match the raw channel scales, so the model spends capacity on scale instead of
shape, and the cumulative-sum integration turns the leftover per-step error into
endpoint overshoot.}
\label{fig:syn-norm}
\end{figure}

This is the channel-scale mismatch from Section~\ref{sec:syn-norm}. The diffusion
prior is isotropic $\mathcal{N}(0,I)$, but the raw $x$-velocity has a much larger
scale than the raw $y$-velocity, so one shared noise level is too large for one
channel and too small for the other. The model ends up modelling the scale rather
than the shape, and the Euler-integrated cumulative sum of
Equation~\eqref{eq:syn-cumsum} turns the leftover per-step errors into endpoint
overshoot. Per-channel normalization costs two 2-vectors, inverts exactly, and
removes the mismatch, so there is little reason not to use it. That is the
empirical argument for keeping it in the production pipeline, including the
online exponential-moving-average variant mentioned in
Section~\ref{sec:syn-norm}.

\begin{table}[htbp]
\centering
\caption{Qualitative summary of the synthetic ablations. All runs use four
temporal blocks and $5000$ training steps. Coverage means how well the generated
set spans the ground-truth fan.}
\label{tab:syn-ablation}
\begin{tblr}{
  colspec={llll},
  row{1}={font=\bfseries},
}
\toprule
Setting & Path smoothness & Endpoint cloud & Scale errors \\
\midrule
$D=8$, norm   & minor wobble & moderate spread on arc & none \\
$D=16$, norm  & smooth       & tight on arc           & none \\
\midrule
with norm     & smooth       & arc-like, $x\lesssim1.0$ & none \\
without norm  & wavy, kinked & diffuse, $x$ up to $\sim1.2$ & overshoot \\
\bottomrule
\end{tblr}
\end{table}

\subsection{Summary}
\label{sec:syn-summary}

The synthetic study checks the generative core of the trajectory probe on a
problem whose answer we know in closed form. The data are cubic B\'ezier curves
running from the origin to a point on a $60^\circ$ unit arc, parameterized by a
turn scalar and represented as zero-padded first-difference velocities. Those
velocities are standardized per channel with cached statistics and inverted
exactly at sampling time. The denoiser is a pointwise MLP encoder feeding a
four-layer DiT-style temporal self-attention stack with AdaLN-Zero conditioning
on the diffusion time, and a linear head emits a per-step flow-matching velocity
field. Training is conditional flow matching on a linear schedule with an MSE
objective against the field $x-(1-\sigma_{\min})\varepsilon$, and sampling is
$20$-step Euler integration from noise to data followed by a cumulative sum into
positions. Since the turn conditioning is switched off, what gets measured is
distribution coverage and not per-turn accuracy. The two ablations back up the
design. A wider hidden size gives a cleaner and better-spread fan with tighter
endpoints, and per-channel velocity normalization turns out to be necessary:
without it the trajectories are wavy, they overshoot, and their endpoints
scatter. So the flow-matching temporal-transformer recipe does learn smooth,
well-distributed trajectories on a controlled problem, which is what we wanted to
establish before adding image conditioning on real nuPlan data with the same
architecture.

    \section{The Probe on Real Driving Logs}
\label{sec:navsim-probe}

The synthetic study in Section~\ref{sec:synthetic} answered a narrow question,
namely whether a flow-matching temporal transformer can denoise a whole
trajectory at once and reproduce a target distribution, under conditions where
nothing else could go wrong: analytic data, zeroed conditioning, and every
episode the same length. This section describes what the system looks like once
those three simplifications are removed and the same generative core is attached
to the frozen Orbis latent on real NAVSIM driving logs. This is the part the
earlier sections put off, and it is where most of the engineering actually sits.

Three properties of real logs drive most of what follows. The conditioning is no
longer a scalar but a \emph{spatial grid of tokenizer latents}, one grid per
observed frame, and it has to be reduced to something a trajectory transformer
can take in without swamping the two-channel velocity signal. Episodes have
\emph{different lengths}, since a NAVSIM scene is only as long as the run of
frames whose camera images happen to be on disk, so batches are ragged. And the
behaviour we care about, turning or branching at an intersection, is rare, so the
pipeline has to keep whole episodes instead of short fixed windows, and the
evaluation has to be able to pick out the episodes where a distribution over
futures means anything at all. Figure~\ref{fig:navsim-pipeline} shows the
pipeline end to end.

\begin{figure}[htbp]
\centering
\includegraphics[width=\textwidth,height=0.38\textheight,keepaspectratio]{../diagrams/diagram-navsim-pipeline.pdf}
\caption{The NAVSIM probing pipeline. Everything above the dashed arrow runs
once: the driving logs are scanned into an episode index, each episode is decoded
into front-camera frames at $2$\,Hz, and the frozen stage-1 tokenizer turns those
frames into per-frame latents written to one HDF5 file per episode. Everything
below is the training loop. A window dataset reads a whole episode's trajectory
but only the leading image latents, a ring buffer hides the HDF5 read latency, a
padding collate assembles ragged episodes into a rectangular batch with explicit
lengths, and the flow-matching probe is trained and evaluated on that batch with
the padded frames masked out.}
\label{fig:navsim-pipeline}
\end{figure}

\subsection{From driving logs to episodes}
\label{sec:navsim-episodes}

Two source datasets sit behind a single configuration switch. They differ only in
how an \emph{episode} is enumerated; the per-frame dictionary is the same in
both, so all the geometry code is shared. The \textbf{navhard} split
(\texttt{navhard\_two\_stage}) gives one episode per meta pickle, taken as its
longest scene, which is around five frames. That is enough for a sanity check and
far too short for a trajectory distribution. The \textbf{navtrain} split (the
\texttt{trainval} logs plus the \texttt{trainval} sensor blobs) is the one the
experiments use: each log pickle spans many scenes of roughly forty frames at
$2$\,Hz, and one episode is one scene.

The complication with navtrain is that the sensor download is usually
\emph{partial}. The log pickles list every frame, but only some of the
corresponding front-camera JPEGs are actually on disk. Instead of dropping whole
scenes or quietly returning black images, we define an episode as the
\emph{longest contiguous run of frames within a scene whose front-camera image is
present}. This keeps the temporal spacing uniform, which matters because a
velocity is a finite difference and a hole in the frame sequence would silently
double one step. The cost is that some episodes get shorter. Scenes with fewer
than \texttt{min\_frames}$\,=5$ usable frames are thrown away. The scan touches
the whole sensor tree, so the resulting (log, scene) $\to$ present-frames mapping
is cached to \texttt{data/navtrain\_episode\_index.pkl} and only rebuilt when
that file is deleted, for instance after downloading more blobs. The cache used
for the runs here holds $9\,577$ episodes.

Each frame contributes a global ego pose: the translation $(x,y)$ from
\texttt{ego2global\_translation}, and a yaw recovered from the
\texttt{ego2global\_rotation} quaternion $[w,x,y,z]$ as
\begin{equation}
\psi = \operatorname{atan2}\!\big(2(wz+xy),\; 1-2(y^2+z^2)\big).
\label{eq:yaw}
\end{equation}
The whole episode is then expressed in the local frame of its \emph{first} frame,
by subtracting the initial position and rotating by $-\psi_0$,
\begin{equation}
p_t \;=\; R(-\psi_0)\,\big(g_t - g_0\big),
\qquad
R(-\psi_0)=\begin{pmatrix}\cos\psi_0 & \sin\psi_0\\ -\sin\psi_0 & \cos\psi_0\end{pmatrix},
\label{eq:local-frame}
\end{equation}
so every trajectory starts at the origin heading along $+x$. The synthetic
B\'ezier curves satisfied the same convention by construction, which is why the
same normalization and the same plotting code can be reused unchanged.

The model again works in velocity space, with one representational difference
that is easy to misread when comparing the two code paths, so it is worth being
explicit. The synthetic data used a \emph{forward} difference with the last entry
zero-padded, while the NAVSIM loader uses a \emph{backward} difference with the
first entry zeroed,
\begin{equation}
v_0 = 0,\qquad v_t = p_t - p_{t-1}\ \ (t\ge 1).
\label{eq:navsim-velocity}
\end{equation}
The advantage is that the inverse is an unshifted cumulative sum,
$\sum_{i\le t} v_i = p_t - p_0$, so integrating the sampled velocities and
re-anchoring at the first frame recovers the positions exactly with no
off-by-one correction. It also means the first predicted entry is structurally
zero, so the model spends one output slot on a constant. That is harmless, but it
is worth remembering when reading per-frame error curves.

The train/validation split is a deterministic function of the scene token rather
than of the file order. A token is hashed with MD5, reduced modulo $1000$, and
assigned to validation if the bucket comes out below
$1000\cdot\texttt{val\_fraction}$, with $\texttt{val\_fraction}=0.1$. Since the
hash depends only on the token, the split survives re-caching, re-indexing and
any change in the number of downloaded blobs, and the qualitative figures use the
same episodes as the logged validation metrics.

\subsection{Caching the frozen tokenizer latents}
\label{sec:navsim-cache}

Encoding images with the stage-1 tokenizer inside the training loop is possible.
A configuration flag (\texttt{encode\_images}) selects it, and the tokenizer is
then constructed lazily on the correct device while the batch carries raw images.
But it is wasteful, because the encoder is frozen and its output for a given
frame never changes. The default path encodes every episode once, in chunks of
$100$ frames, and writes one HDF5 file per episode holding both tokenizer
streams, \texttt{encoded\_q\_rec} and \texttt{encoded\_q\_sem}, each of shape
$(T, C, H_{\text{lat}}, W_{\text{lat}})$, along with the trajectory, the
velocity, the frame rate and the scene token. With the $\times16$ spatial
compression from Section~\ref{sec:orbis} and the $256\times256$ centre crop the
loader uses, the latent grid is $16\times16$ positions with $C=16$ channels.

The probe consumes the \emph{semantic} stream \texttt{encoded\_q\_sem}. The
semantic encoder is the one regularized towards DINOv2 features, so it is the
stream that should carry scene structure such as lanes, vehicles and free space,
whereas the detail stream is pushed towards the low-level appearance needed for
faithful reconstruction. We cache the reconstruction stream anyway, since it
costs almost nothing extra, so the choice can be ablated later without
re-encoding the dataset.

\subsection{Whole episodes, ragged batches, and the frame mask}
\label{sec:navsim-padding}

The first version of the training loop cut a fixed window of
\texttt{num\_frames} frames out of each episode. That made every batch
rectangular and every tensor shape static, and it also threw away most of the
data. With episodes of about forty frames and a window of eight, the model never
saw a trajectory long enough to contain a complete manoeuvre, and the longest and
most interesting episodes were truncated just like the shortest ones. Setting
\texttt{num\_frames}$\,=0$ switches the dataset to serving \emph{whole}
episodes, at which point the batch becomes ragged.

Ragged batches are handled at the collate level, not inside the model. A custom
collate function pads every time-major tensor of the sample dictionary with zeros
up to the longest episode in the batch, using \texttt{pad\_sequence} so that
trailing dimensions are left alone, and records the true lengths in
\texttt{batch["metadata"]["length"]}. Padding is applied per key independently.
That matters because two different time axes live in one sample: the trajectory
and velocity have length $T$, while the image latents have length
$T_{\text{img}}\le\texttt{context\_images}$, and a very short episode can end up
with fewer image frames than the nominal image context.

The model turns the recorded lengths into a boolean frame mask
$m\in\{0,1\}^{B\times T}$ with $m_{b,t}=1$ iff $t<\ell_b$. The mask is used for
\emph{selection} and not for truncation, which is the part that is easy to get
wrong: the network still predicts all $T_{\max}$ frames of every sample, and the
mask is what removes the padded tail from the loss, from the online velocity
statistics and from every validation metric. Writing the flow-matching residual
for episode $b$ and frame $t$ as $r_{b,t}=\hat u_{b,t}-u_{b,t}\in\mathbb{R}^2$,
the objective is a masked mean over the frames that are both predicted (past the
trajectory context $c$) and valid,
\begin{equation}
\mathcal{L}
= \frac{\displaystyle\sum_{b}\sum_{t\ge c} m_{b,t}\,\tfrac{1}{2}\lVert r_{b,t}\rVert_2^2}
       {\displaystyle\sum_{b}\sum_{t\ge c} m_{b,t}},
\label{eq:masked-loss}
\end{equation}
This is implemented without a per-sample loop, so variable lengths cost one
elementwise multiply. Figure~\ref{fig:navsim-mask} shows the layout of a padded
batch.

\begin{figure}[htbp]
\centering
\includegraphics[width=0.86\textwidth,height=0.34\textheight,keepaspectratio]{../diagrams/diagram-navsim-mask.pdf}
\caption{A padded batch of three episodes. Every episode is stored at the full
padded length $T_{\max}=8$ and the network predicts all $T_{\max}$ velocity steps
for all three. The recorded lengths $\ell_b$ become the frame mask $m_{b,t}$,
which removes the padded tail from the loss of
Equation~\eqref{eq:masked-loss}, from the exponential-moving-average velocity
statistics and from the validation metrics. Only the leading
\texttt{context\_images} frames, here two, carry an image latent, and those are
the only latents read from disk.}
\label{fig:navsim-mask}
\end{figure}

Reading whole episodes would multiply the amount of latent data pulled off disk
if it were not for the second half of the change. The dataset slices the HDF5
arrays so that \emph{only} the leading \texttt{context\_images} latents are read,
while the trajectory, which is two floats per frame, is always read in full. The
denoiser never looks at an image beyond the context, so nothing is lost, and the
I/O per sample becomes essentially independent of episode length. What latency is
left gets absorbed by an in-memory ring buffer, which serves repeated samples and
pulls \texttt{replace\_count}$\,=10$ fresh episodes from the source dataset every
$40$ served samples, growing to at most $20\,000$ cached episodes. The buffer
trades a controlled amount of sample staleness for a GPU that is never starved by
HDF5 decompression.

\subsection{What changed in the image aggregator}
\label{sec:navsim-aggregator}

The conditioning path has been revised more than anything else, and that revision
is the main architectural difference between the model described in
Section~\ref{sec:probe} and the one we actually train.

\paragraph{The original design: one pooled token per frame, concatenated.}
In the first version the aggregator reduced each frame's latent grid to a
\emph{single} token. The $16$ tokenizer channels were lifted to the model width
$D$ by a linear map, the $H_{\text{lat}}\!\times\!W_{\text{lat}}$ positions were
flattened into a sequence with a 2D sinusoidal position embedding, and one
cross-attention query pooled that sequence into a single $D$-dimensional vector,
with the frame embedding entering as the AdaLN conditioning of the pooling block.
The result was a tensor of shape $(B,T,D)$, one image vector per frame, whose
entries beyond the image context were explicitly zeroed. That vector was then
\emph{concatenated to the per-point feature} of every trajectory point, so the
input to the pointwise MLP had width $2+3D$: two velocity channels, the
diffusion-time embedding, the frame embedding and the frame's image token.

Three problems with this only show up on real data. The capacity bottleneck is
severe, since an entire street scene is squeezed into $D=64$ numbers before the
trajectory model sees any of it, so the probe has no way to ask different
questions of different parts of the image. The alignment is rigid: because the
image token is concatenated per frame, image information reaches trajectory point
$t$ only if frame $t$ carried an image, which forces the awkward zeroing of all
non-context frames and burns an input slot on a constant zero vector for most of
the episode. And the interface is brittle, because the width of the first linear
layer depends on how image evidence is supplied, so changing the number of
observed frames changes the parameter shapes.

\paragraph{The current design: several tokens per frame, as attention keys.}
The aggregator now emits $N=\texttt{image\_tokens}$ tokens per observed frame
instead of one. The same cross-attention block is used with $N$ queries rather
than one, giving $(B,T_{\text{img}},N,D)$, and a 1D sinusoidal position embedding
over the $N$ slots is added so the slots stay distinguishable once they are
flattened together with the frames. Only the leading
$T_{\text{img}}=\texttt{context\_images}$ frames get aggregated at all, so the
zeroing of later frames goes away; later frames are simply never encoded. The
resulting tokens are flattened into one sequence of length
$T_{\text{img}}\!\cdot\!N$.

The important part is that these tokens no longer enter the per-point features.
They are prepended to the key/value sequence of every temporal block, so the
per-point input shrinks to width $2+2D$ and the image evidence reaches the
trajectory through \emph{cross-attention} rather than concatenation. Each
predicted frame can now attend to whichever image tokens it needs, the number of
observed frames and the number of tokens per frame can both change without
touching a parameter shape, and the conditioning bandwidth is set by
$T_{\text{img}}\cdot N$ instead of by the model width. In the configuration used
here, $\texttt{context\_images}=2$ and $\texttt{image\_tokens}=9$ give $18$ image
tokens per episode. Figure~\ref{fig:navsim-aggregator} contrasts the two
designs.

There is one asymmetry in the current implementation that should be recorded.
Slot identity enters the image tokens \emph{additively}, through the 1D position
embedding over the $N$ slots, while frame identity enters only
\emph{multiplicatively}, through the AdaLN modulation of the pooling block by the
frame embedding. With $T_{\text{img}}=2$ the two observed frames are distinguished
by a modulation of the pooling rather than by an explicit temporal tag on the
tokens themselves. Adding a frame embedding to the tokens would make the temporal
ordering of the image context explicit, and that is the obvious next refinement
if the image context is ever made longer. One more detail, since
Section~\ref{sec:probe} describes the probe in the abstract: the pooling queries
are \emph{fixed} sinusoidal position embeddings, not learned parameters. What is
learned is the projection of those queries, the key/value projections of the
latent grid, the AdaLN modulation and the feed-forward branch. The query tokens
themselves are deterministic slot identifiers.

\begin{figure}[htbp]
\centering
\begin{subfigure}[t]{0.47\textwidth}
\centering
\includegraphics[width=\linewidth,height=0.32\textheight,keepaspectratio]{../diagrams/diagram-aggregator-old.pdf}
\caption{Before: a single pooled token per frame, concatenated to every
trajectory point.}
\label{fig:navsim-aggregator-old}
\end{subfigure}
\hfill
\begin{subfigure}[t]{0.47\textwidth}
\centering
\includegraphics[width=\linewidth,height=0.32\textheight,keepaspectratio]{../diagrams/diagram-aggregator-new.pdf}
\caption{After: $N$ tokens per observed frame, entering the attention stack as
keys and values.}
\label{fig:navsim-aggregator-new}
\end{subfigure}
\caption{The revision of the image aggregator. Pooling the latent grid into one
vector per frame (left) pushed the whole scene through a $D$-dimensional
bottleneck, tied image evidence to the frame index by concatenation, and required
non-context frames to be zeroed. Emitting $N$ tokens per observed frame and
routing them through cross-attention (right) widens the conditioning bandwidth to
$T_{\text{img}}\cdot N$ tokens, lets every predicted frame query the image context
freely, removes the zero-padding of non-context frames, and makes the per-point
input width independent of how much image evidence is supplied.}
\label{fig:navsim-aggregator}
\end{figure}

\subsection{The denoiser as a whole}
\label{sec:navsim-denoiser}

With the aggregator settled, the rest of the denoiser is a small
sequence-to-sequence transformer whose only trainable inputs are the noisy
velocities and two scalar embeddings. Given the noisy velocity trajectory
$\hat v\in\mathbb{R}^{B\times T\times2}$ at diffusion time $t$, the network forms
a diffusion-time embedding $e_{\text{time}}\in\mathbb{R}^{B\times D}$ broadcast
across frames, and a frame embedding $e_{\text{frame}}\in\mathbb{R}^{T\times D}$
computed from the \emph{normalized} index $\operatorname{arange}(T)/T$ so that it
is well defined for any episode length. Each point is then described by
\begin{equation}
x_i=\big[\underbrace{\hat v_i}_{2},\;\underbrace{e_{\text{time}}}_{D},\;
\underbrace{e_{\text{frame},i}}_{D}\big]\in\mathbb{R}^{2+2D},
\label{eq:navsim-input}
\end{equation}
and split at the trajectory-context boundary: the first
$c=\texttt{context\_traj}$ points go through a \emph{context} MLP and the
remaining $T-c$ through a \emph{prediction} MLP. There are two MLPs rather than
one because the two groups do different jobs. The context tokens are known and
never updated; the prediction tokens are the query stream the stack refines.

Each temporal block is a DiT-style cross-attention block with AdaLN-Zero
conditioning on $e_{\text{time}}$, as in Section~\ref{sec:syn-model}, but here it
is used in a genuinely \emph{cross}-attentive configuration. The query is the
current prediction stream, and the key/value sequence is the concatenation of the
image tokens, the trajectory-context tokens and the prediction tokens,
\begin{equation}
\text{pred}\;\leftarrow\;
\text{TemporalBlock}\big(\underbrace{[\,\text{img}\;\|\;\text{ctx}\;\|\;\text{pred}\,]}_{\text{keys/values}},\;
\underbrace{\text{pred}}_{\text{queries}},\; e_{\text{time}}\big),
\label{eq:navsim-block}
\end{equation}
repeated for every block in the stack. The key sequence is rebuilt from the
updated predictions at each layer, so the image and context tokens stay static
while the prediction stream is refined. That is closer to a decoder stack with a
fixed memory than to plain self-attention. A LayerNorm and a linear map project
the final prediction stream back to two channels, which gives the velocity field
on the predicted frames only, of shape $(B,T-c,2)$.
Figure~\ref{fig:navsim-denoiser} shows the assembled network.
Figure~\ref{fig:navsim-sketch} is a simpler version of the same picture, next
to the tokenizer that produces the image latents.

The configuration used here sets $D=64$ with $8$ attention heads and a
feed-forward width of $2D$, and the stack currently holds two temporal blocks
(reduced from four when the image tokens were added, which lengthened the key
sequence and therefore the cost of each block). With
$\texttt{context\_traj}=0$, $\texttt{context\_images}=2$ and
$\texttt{image\_tokens}=9$, a block attends over $18+T$ keys and produces $T$
queries.

\begin{figure}[htbp]
\centering
\includegraphics[width=\textwidth,height=0.42\textheight,keepaspectratio]{../diagrams/diagram-navsim-denoiser.pdf}
\caption{The trajectory denoiser on real data. The image branch (top) turns the
cached tokenizer latents of the $T_{\text{img}}$ observed frames into
$T_{\text{img}}\cdot N$ tokens. The trajectory branch (bottom left) encodes each
point from the noisy velocity, the diffusion-time embedding and the frame
embedding, splitting at the trajectory-context boundary. Every temporal block
takes the concatenated image, context and prediction tokens as keys and values
and the prediction stream as queries, conditioned on the diffusion time through
AdaLN-Zero, and the refreshed prediction stream re-enters the key sequence at the
next layer. A LayerNorm and a linear head emit the flow-matching velocity field
on the predicted frames only.}
\label{fig:navsim-denoiser}
\end{figure}

\begin{figure}[htbp]
\centering
\begin{subfigure}[t]{0.38\textwidth}
\centering
\includegraphics[width=\linewidth]{../diagrams/diagram-tokenizer.pdf}
\caption{The hybrid tokenizer. The semantic and detail encoders each give an
$18\times32$ grid with $16$ channels, and the two grids are concatenated along
the channels.}
\label{fig:navsim-sketch-tokenizer}
\end{subfigure}
\hfill
\begin{subfigure}[t]{0.58\textwidth}
\centering
\includegraphics[width=\linewidth]{../diagrams/diagram-denoiser-sketch.pdf}
\caption{The denoiser. The spatial block turns each context frame into tokens,
which join the noisy trajectory points in the temporal block. The diffusion
time step comes in through AdaLN-Zero.}
\label{fig:navsim-sketch-denoiser}
\end{subfigure}
\caption{Simplified view of the tokenizer and the denoiser. Only one image
token per frame is drawn, the model uses $N$. The resolution in (a) is the one
the world model was trained at. The probe in this section crops to
$256\times256$ instead, which gives a $16\times16$ grid.}
\label{fig:navsim-sketch}
\end{figure}

\subsection{Two independent contexts}
\label{sec:navsim-contexts}

The model carries \emph{two} unrelated notions of context, and they are easy to
conflate even though the code keeps them separate. The \emph{trajectory context}
$\texttt{context\_traj}=c$ is the number of leading frames whose ground-truth
velocities are handed to the model and therefore not predicted. Those frames are
copied into the noisy trajectory before the network runs, held fixed during ODE
integration, and left out of the loss. The \emph{image context}
$\texttt{context\_images}=T_{\text{img}}$ is the number of leading frames that
carry an image latent. The two are set independently and the episode length $T$
is free, so everything after the trajectory context gets predicted no matter how
long the episode is.

The runs reported here use $c=0$ and $T_{\text{img}}=2$. This is the strictest
setting for a probe, and the most informative one: the model gets two
front-camera observations and \emph{no} ground-truth motion at all, and has to
produce the whole trajectory from the frozen latent. It also keeps the question
from Section~\ref{sec:probe} free of trajectory extrapolation. With $c>0$ a
constant-velocity continuation of the given prefix would already score well, so
$c=0$ is what isolates the image representation. The context branch is still kept
alive in the architecture, because the planned experiments include training with
the context trajectory supplied half the time. With $c=0$ the context MLP just
receives an empty slice and contributes no tokens.

One implementation detail looks like information leakage on a first reading and
is not: the sampler passes the ground-truth velocities to the denoiser at
inference time. They are only used to fill the first $c$ entries of the
trajectory, so with $c=0$ they are never read. The interface is written this way
so that one call site serves both settings.

\subsection{Online velocity normalization}
\label{sec:navsim-norm}

Section~\ref{sec:syn-norm} showed why per-channel velocity standardization
matters: an isotropic diffusion prior does not match two channels of very
different scale, and the mismatch turns into endpoint drift after the cumulative
sum. On real logs the argument is if anything stronger, since the longitudinal
velocity of a car at $2$\,Hz is on the order of metres per frame while the
lateral component is near zero on straights. But the one-shot computation used in
the synthetic study is inconvenient here, because the statistics would have to be
recomputed every time the dataset, the split or the frame rate changed.

The production path therefore keeps the same transform and estimates its
parameters online. A small module maintains an exponential moving average of the
per-channel mean and variance with decay $0.999$, seeded directly by the first
batch it sees,
\begin{equation}
\mu \leftarrow \gamma\,\mu + (1-\gamma)\,\mu_{\text{batch}},
\qquad
\sigma^2 \leftarrow \gamma\,\sigma^2 + (1-\gamma)\,\sigma^2_{\text{batch}},
\qquad \gamma = 0.999,
\label{eq:ema}
\end{equation}
and applies the familiar pair
$\hat v=(v-\mu)/(\sigma+\varepsilon)$ and $v=\hat v(\sigma+\varepsilon)+\mu$.
Three details make this work here. The batch statistics are computed over the
\emph{valid} frames only, using the same mask as the loss. Without that, the
zero-padded tail of every short episode would drag the mean and the variance
towards zero, and the more ragged the batch the worse the distortion. The
statistics are stored as non-trainable parameters, so they are part of the
checkpoint and a reloaded model denormalizes the same way the trained one did.
That is easy to get wrong when the statistics live in the dataset instead. And
the update only happens while the module is in training mode, so validation and
sampling read the statistics without disturbing them.

\subsection{Training, sampling and configuration}
\label{sec:navsim-training}

The objective is the same as in Section~\ref{sec:syn-flow}: a linear
flow-matching schedule with $\alpha(t)=1-t$ and
$\sigma(t)=\sigma_{\min}+t(1-\sigma_{\min})$, one diffusion time drawn uniformly
per episode, and a mean-squared regression onto the velocity field
$u=x-(1-\sigma_{\min})\varepsilon$. Two things differ. The average is the masked
one of Equation~\eqref{eq:masked-loss}, and only the frames past the trajectory
context contribute.

Sampling mirrors the synthetic procedure with the context handling added. The
trajectory is initialized from pure noise, the first $c$ entries are overwritten
by the known context velocities, and the ODE is integrated backwards from $t=1$
to $t=0$ with $50$ uniform Euler steps, updating only the predicted suffix at
each step. The integrated velocities are denormalized with the EMA statistics and
turned into positions by a cumulative sum re-anchored at the first frame. By
Equation~\eqref{eq:navsim-velocity} that places the predicted path at the origin
in the ego frame of the first observed frame, which is the same frame the cached
trajectory and the map BEV live in, so predictions, ground truth and map overlay
with no further transform.

The validation step logs three quantities of different character on purpose.
\texttt{val/loss} is the \emph{same} masked flow-matching objective as
\texttt{train/loss} evaluated on held-out episodes, so the two curves can be read
against each other and the generalization gap is visible at a glance. It costs
one forward pass. \texttt{val/mse} and \texttt{val/std} need full ODE rollouts
instead, so they measure the sampled trajectories rather than the regression
target. Table~\ref{tab:navsim-config} lists the configuration.

\begin{table}[htbp]
\centering
\caption{Configuration of the NAVSIM probe. The values are those of
\texttt{exp\_navsim/config.yaml} for the runs discussed here. The episode length
$T$ is not a hyperparameter; it is a property of each batch.}
\label{tab:navsim-config}
\begin{tblr}{
  colspec={llX[l]},
  row{1}={font=\bfseries},
}
\toprule
Parameter & Value & Role \\
\midrule
source dataset        & navtrain     & full NAVSIM training logs, $\sim\!40$-frame scenes \\
frame rate            & $2$\,Hz      & stored rate; no subsampling \\
episodes (cached)     & $9\,577$     & longest on-disk contiguous run per scene \\
image size            & $256\times256$ & resize $+$ centre crop of \texttt{CAM\_F0} \\
latent stream         & \texttt{encoded\_q\_sem} & DINOv2-regularized semantic latents \\
latent channels $C$   & $16$         & tokenizer output width \\
\midrule
$\texttt{context\_traj}$   & $0$     & no ground-truth motion given; all frames predicted \\
$\texttt{context\_images}$ & $2$     & observed frames carrying a latent \\
$\texttt{image\_tokens}$   & $9$     & tokens the spatial block emits per observed frame \\
hidden width $D$      & $64$         & embeddings, attention and feed-forward width \\
temporal blocks       & $2$          & cross-attention denoiser stack \\
attention heads       & $8$          & per block, dropout $0.1$ \\
\midrule
$\sigma_{\min}$       & $10^{-5}$    & flow-matching schedule floor \\
diffusion steps       & $50$         & Euler steps at sampling time \\
EMA decay             & $0.999$      & online velocity statistics \\
batch size            & $100$        & episodes per step, zero-padded to $T_{\max}$ \\
base LR               & $5\cdot10^{-4}$ & AdamW, weight decay $0.01$, cosine decay to $0.1\times$ \\
warmup                & $0$ steps    & \\
precision             & 16-mixed     & \\
validation samples    & $20$         & rollouts per context, for \texttt{val/std} \\
validation interval   & every $5$ epochs & \\
\bottomrule
\end{tblr}
\end{table}

\subsection{Metrics for a distribution of futures}
\label{sec:navsim-metrics}

A probe that is asked for a \emph{distribution} cannot be judged by pointwise
error alone, because the two failure modes we care about sit at opposite ends of
the same axis. A model that collapses onto the mean of the data will look
respectable under a squared error and be useless as a world-model read-out. A
model that samples wildly will cover every future and be just as useless. So we
log two quantities.

The first is the masked trajectory MSE, averaged over the predicted and valid
frames of all sampled rollouts, which says how close the samples sit to the one
future that actually happened. The second measures spread. The \emph{total turn}
of a trajectory is the accumulated change of heading along it: with segments
$s_t=p_{t+1}-p_t$ and headings $\theta_t=\operatorname{atan2}(s_{t,y},s_{t,x})$,
\begin{equation}
\Theta \;=\; \sum_{t} \operatorname{wrap}\!\big(\theta_{t+1}-\theta_t\big),
\qquad
\operatorname{wrap}(\delta)=\operatorname{atan2}(\sin\delta,\cos\delta),
\label{eq:total-turn}
\end{equation}
The wrapping keeps the increments in $[-\pi,\pi]$, so a heading crossing the
$\pm\pi$ branch cut does not register as a full turn. A turn increment at index
$t$ is built from two consecutive segments and therefore from three consecutive
frames, so an increment is only counted when the third of those frames is valid.
The padded tail contributes nothing, which is what keeps the metric comparable
across episodes of different length. Sampling the same context $N=20$ times,
taking the standard deviation of $\Theta$ across those samples, and averaging
over the batch,
\begin{equation}
\texttt{val/std} \;=\; \frac{1}{B}\sum_{b} \operatorname{std}_{n\le N}\ \Theta_{b}^{(n)},
\label{eq:turn-std}
\end{equation}
gives an interpretable read-out of how much the probe disagrees with itself about
where the car is going. The pair has to be read together. A low MSE with a
near-zero turn spread is the signature of mode collapse onto the average path,
while a growing spread that does not hurt the MSE is evidence of real
multi-modality.

The qualitative counterpart is a per-episode PDF, produced by a separate
evaluation script from the same deterministic validation split. Each page shows
the first and last context front views, the NAVSIM map bird's-eye view with the
\emph{whole} sampled distribution drawn on it, and the trajectory panels. The
model predicts the full padded length for every sample, so the page trims each
episode back to its true length before plotting. Episodes are pre-filtered by
driven distance ($\ge 5$\,m) and by the angle between the initial heading and the
start-to-end vector ($\ge 10^\circ$). A stationary or perfectly straight episode
makes any trajectory distribution look correct, so the filter is there to pick
out the cases where the question is not trivial.

\subsection{Status, limitations and planned experiments}
\label{sec:navsim-status}

The pipeline above is the state the system currently runs in: whole-length
episodes, padded batches with a masked objective, the multi-token image
aggregator feeding a cross-attention denoiser, online velocity normalization, and
a validation protocol reporting both accuracy and spread. Going from fixed short
windows to full episodes, together with the aggregator revision, is what moved
the qualitative predictions from unconvincing to clearly structured.

Several limitations are known and should be stated plainly. The attention stack
does not yet get a key-padding mask, so padded trajectory tokens are visible as
keys even though the loss excludes them. Those tokens are zero-valued and the
effect looks small, but it is a real inconsistency between the loss and the
attention, and it is the next thing to fix. The image context is short
($T_{\text{img}}=2$) and its frame identity enters only through AdaLN modulation,
as discussed in Section~\ref{sec:navsim-aggregator}. The probe is evaluated with
MSE and turn spread rather than with the displacement-based metrics the driving
literature uses, and there is still no control baseline, meaning a
constant-velocity extrapolation or a probe trained on a shuffled latent, to
establish how much of the performance comes from the representation rather than
from the strong prior structure of driving trajectories. Finally, and this
matters most for the claim the project set out to test, everything reported here
conditions on \emph{stage-1 tokenizer} latents. The stage-2 world model of
Section~\ref{sec:orbis} has not been probed at all.

That last gap is what the planned continuation is about, and it has three parts.
The first is a layer-wise probe of the world model: freeze the STDiT stack, tap
its intermediate block activations, and train the same probe on each tap in turn
to find where ego-motion information is strongest as a function of depth. A flat
depth curve and a peaked one would both tell us something. The second is a
capacity and control sweep that fixes the best tap and varies probe size against
baselines, so that a good score can be credited to the representation rather than
to the read-out. The third goes at calibration directly: sample
$k\in\{1,5,10,20\}$ futures per context on strata separating free-flow driving
from branching intersections, and report $\text{minADE}_k$ and mode recall on map
branches next to a sharpness measure on straights, so a model cannot score well
just by sampling everything. That third experiment needs longer and denser source
clips, since a five-second NAVSIM scene rarely contains a real branch point, and
the data ingest for that is the one piece of infrastructure still missing.


    \begin{thebibliography}{99}

    \bibitem{orbis}
    A.~Mousakhan, S.~Mittal, S.~Galesso, K.~Farid, and T.~Brox.
    \newblock Orbis: Overcoming challenges of long-horizon prediction in driving
    world models.
    \newblock \emph{arXiv preprint arXiv:2507.13162}, 2025.

    \bibitem{ho_ddpm}
    J.~Ho, A.~Jain, and P.~Abbeel.
    \newblock Denoising diffusion probabilistic models.
    \newblock In \emph{NeurIPS}, 2020.

    \bibitem{lipman_fm}
    Y.~Lipman, R.~T.~Q. Chen, H.~Ben-Hamu, M.~Nickel, and M.~Le.
    \newblock Flow matching for generative modeling.
    \newblock In \emph{ICLR}, 2023.

    \bibitem{dit}
    W.~Peebles and S.~Xie.
    \newblock Scalable diffusion models with transformers.
    \newblock In \emph{ICCV}, 2023.

    \bibitem{nwm}
    A.~Bar, G.~Zhou, D.~Tran, T.~Darrell, and Y.~LeCun.
    \newblock Navigation world models.
    \newblock \emph{arXiv preprint arXiv:2412.03572}, 2024.

    \bibitem{maskgit}
    H.~Chang, H.~Zhang, L.~Jiang, C.~Liu, and W.~T. Freeman.
    \newblock MaskGIT: Masked generative image transformer.
    \newblock In \emph{CVPR}, 2022.

    \bibitem{vqgan}
    P.~Esser, R.~Rombach, and B.~Ommer.
    \newblock Taming transformers for high-resolution image synthesis.
    \newblock In \emph{CVPR}, 2021.

    \bibitem{dosovitskiy_vit}
    A.~Dosovitskiy, L.~Beyer, A.~Kolesnikov, et~al.
    \newblock An image is worth 16x16 words: Transformers for image recognition
    at scale.
    \newblock In \emph{ICLR}, 2021.

    \bibitem{nuplan}
    H.~Caesar, J.~Kabzan, K.~S. Tan, et~al.
    \newblock nuPlan: A closed-loop ML-based planning benchmark for autonomous
    vehicles.
    \newblock \emph{arXiv preprint arXiv:2106.11810}, 2021.

    \bibitem{navsim}
    W.~Cao, M.~Hallgarten, T.~Li, et~al.
    \newblock Pseudo-simulation for autonomous driving.
    \newblock In \emph{CoRL}, 2025.

    \bibitem{hydra}
    K.~Li, Z.~Li, S.~Lan, et~al.
    \newblock Hydra-MDP++: Advancing end-to-end driving via expert-guided
    Hydra-distillation.
    \newblock \emph{arXiv preprint arXiv:2503.12820}, 2025.

    % ---------- Planning literature (Section 2.4) ----------

    \bibitem{e2esurvey}
    L.~Chen, P.~Wu, K.~Chitta, B.~Jaeger, A.~Geiger, and H.~Li.
    \newblock End-to-end autonomous driving: Challenges and frontiers.
    \newblock \emph{IEEE TPAMI}, 46(12):10164--10183, 2024.

    \bibitem{e2eplansurvey}
    Planning-oriented end-to-end autonomous driving: Architectures, evaluation,
    and emerging paradigms.
    \newblock \emph{arXiv preprint arXiv:2608.20111}, 2026.

    \bibitem{idm}
    M.~Treiber, A.~Hennecke, and D.~Helbing.
    \newblock Congested traffic states in empirical observations and microscopic
    simulations.
    \newblock \emph{Physical Review E}, 62(2):1805--1824, 2000.

    \bibitem{pdm}
    D.~Dauner, M.~Hallgarten, A.~Geiger, and K.~Chitta.
    \newblock Parting with misconceptions about learning-based vehicle motion
    planning.
    \newblock In \emph{CoRL}, 2023.

    \bibitem{alvinn}
    D.~A. Pomerleau.
    \newblock ALVINN: An autonomous land vehicle in a neural network.
    \newblock In \emph{NeurIPS}, 1988.

    \bibitem{pilotnet}
    M.~Bojarski, D.~Del~Testa, D.~Dworakowski, et~al.
    \newblock End to end learning for self-driving cars.
    \newblock \emph{arXiv preprint arXiv:1604.07316}, 2016.

    \bibitem{cil}
    F.~Codevilla, M.~M\"{u}ller, A.~L\'{o}pez, V.~Koltun, and A.~Dosovitskiy.
    \newblock End-to-end driving via conditional imitation learning.
    \newblock In \emph{ICRA}, 2018.

    \bibitem{chauffeurnet}
    M.~Bansal, A.~Krizhevsky, and A.~Ogale.
    \newblock ChauffeurNet: Learning to drive by imitating the best and
    synthesizing the worst.
    \newblock In \emph{RSS}, 2019.

    \bibitem{causalconfusion}
    P.~de~Haan, D.~Jayaraman, and S.~Levine.
    \newblock Causal confusion in imitation learning.
    \newblock In \emph{NeurIPS}, 2019.

    \bibitem{lbc}
    D.~Chen, B.~Zhou, V.~Koltun, and P.~Kr\"{a}henb\"{u}hl.
    \newblock Learning by cheating.
    \newblock In \emph{CoRL}, 2020.

    \bibitem{urbandriver}
    O.~Scheel, L.~Bergamini, M.~Wolczyk, B.~Osi\'{n}ski, and P.~Ondruska.
    \newblock Urban Driver: Learning to drive from real-world demonstrations
    using policy gradients.
    \newblock In \emph{CoRL}, 2021.

    \bibitem{transfuser}
    K.~Chitta, A.~Prakash, B.~Jaeger, Z.~Yu, K.~Renz, and A.~Geiger.
    \newblock TransFuser: Imitation with transformer-based sensor fusion for
    autonomous driving.
    \newblock \emph{IEEE TPAMI}, 45(11):12878--12895, 2023.

    \bibitem{tcp}
    P.~Wu, X.~Jia, L.~Chen, J.~Yan, H.~Li, and Y.~Qiao.
    \newblock Trajectory-guided control prediction for end-to-end autonomous
    driving.
    \newblock In \emph{NeurIPS}, 2022.

    \bibitem{uniad}
    Y.~Hu, J.~Yang, L.~Chen, et~al.
    \newblock Planning-oriented autonomous driving.
    \newblock In \emph{CVPR}, 2023.

    \bibitem{vad}
    B.~Jiang, S.~Chen, Q.~Xu, et~al.
    \newblock VAD: Vectorized scene representation for efficient autonomous
    driving.
    \newblock In \emph{ICCV}, 2023.

    \bibitem{vadv2}
    S.~Chen, B.~Jiang, H.~Gao, et~al.
    \newblock VADv2: End-to-end vectorized autonomous driving via probabilistic
    planning.
    \newblock \emph{arXiv preprint arXiv:2402.13243}, 2024.

    \bibitem{paradrive}
    X.~Weng, B.~Ivanovic, Y.~Wang, Y.~Wang, and M.~Pavone.
    \newblock PARA-Drive: Parallelized architecture for real-time autonomous
    driving.
    \newblock In \emph{CVPR}, 2024.

    \bibitem{hydra_mdp}
    Z.~Li, K.~Li, S.~Wang, et~al.
    \newblock Hydra-MDP: End-to-end multimodal planning with multi-target
    hydra-distillation.
    \newblock \emph{arXiv preprint arXiv:2406.06978}, 2024.

    \bibitem{diffusiondrive}
    B.~Liao, S.~Chen, H.~Yin, et~al.
    \newblock DiffusionDrive: Truncated diffusion model for end-to-end
    autonomous driving.
    \newblock In \emph{CVPR}, 2025.

    \bibitem{diffusionplanner}
    Y.~Zheng, R.~Liang, K.~Zheng, et~al.
    \newblock Diffusion-based planning for autonomous driving with flexible
    guidance.
    \newblock In \emph{ICLR}, 2025.

    \bibitem{gtrs}
    Z.~Li, S.~Lan, et~al.
    \newblock Generalized trajectory scoring for end-to-end multimodal planning.
    \newblock \emph{arXiv preprint arXiv:2506.06664}, 2025.

    \bibitem{drivesuprim}
    DriveSuprim: Towards precise trajectory selection for end-to-end planning.
    \newblock \emph{arXiv preprint arXiv:2506.06659}, 2025.

    \bibitem{carl}
    B.~Jaeger, D.~Dauner, J.~Beau{\ss}wein, S.~Gerstenecker, K.~Chitta, and
    A.~Geiger.
    \newblock CaRL: Learning scalable planning policies with simple rewards.
    \newblock In \emph{CoRL}, 2025.

    \bibitem{think2drive}
    Q.~Li, X.~Jia, S.~Wang, and J.~Yan.
    \newblock Think2Drive: Efficient reinforcement learning by thinking with
    latent world model for autonomous driving.
    \newblock In \emph{ECCV}, 2024.

    \bibitem{emma}
    J.-J. Hwang, R.~Xu, H.~Lin, et~al.
    \newblock EMMA: End-to-end multimodal model for autonomous driving.
    \newblock \emph{arXiv preprint arXiv:2410.23262}, 2024.

    \bibitem{drivevlm}
    X.~Tian, J.~Gu, B.~Li, et~al.
    \newblock DriveVLM: The convergence of autonomous driving and large
    vision-language models.
    \newblock In \emph{CoRL}, 2024.

    \bibitem{senna}
    B.~Jiang, S.~Chen, B.~Liao, et~al.
    \newblock Senna: Bridging large vision-language models and end-to-end
    autonomous driving.
    \newblock \emph{arXiv preprint arXiv:2410.22313}, 2024.

    \bibitem{drivewm}
    Y.~Wang, J.~He, L.~Fan, H.~Li, Y.~Chen, and Z.~Zhang.
    \newblock Driving into the future: Multiview visual forecasting and planning
    with world model for autonomous driving.
    \newblock In \emph{CVPR}, 2024.

    \bibitem{wote}
    Y.~Li, L.~Chen, et~al.
    \newblock End-to-end driving with online trajectory evaluation via BEV world
    model.
    \newblock In \emph{ICCV}, 2025.

    \bibitem{law}
    Y.~Li, L.~Fan, J.~He, et~al.
    \newblock Enhancing end-to-end autonomous driving with latent world model.
    \newblock In \emph{ICLR}, 2025.

    \bibitem{world4drive}
    Y.~Zheng, et~al.
    \newblock World4Drive: End-to-end autonomous driving via intention-aware
    physical latent world model.
    \newblock In \emph{ICCV}, 2025.

    \bibitem{ppgeo}
    P.~Wu, L.~Chen, H.~Li, X.~Jia, J.~Yan, and Y.~Qiao.
    \newblock Policy pre-training for autonomous driving via self-supervised
    geometric modeling.
    \newblock In \emph{ICLR}, 2023.

    \bibitem{vidar}
    Z.~Yang, L.~Chen, Y.~Sun, and H.~Li.
    \newblock Visual point cloud forecasting enables scalable autonomous driving.
    \newblock In \emph{CVPR}, 2024.

    \bibitem{admlp}
    J.-T. Zhai, Z.~Feng, J.~Du, et~al.
    \newblock Rethinking the open-loop evaluation of end-to-end autonomous
    driving in nuScenes.
    \newblock \emph{arXiv preprint arXiv:2305.10430}, 2023.

    \bibitem{egostatus}
    Z.~Li, Z.~Yu, S.~Lan, et~al.
    \newblock Is ego status all you need for open-loop end-to-end autonomous
    driving?
    \newblock In \emph{CVPR}, 2024.

    \end{thebibliography}

    % ---------- End of Document ----------

\end{document}
